\chapter{青春期}\label{ch:g8:puberty}

大约在九岁至十六岁之间，每个人的身体都会
重新塑造自己：第二次生长高峰（见
\cref{rem:g3:stages-of-life:bursts}），远不止身高的变化。
\Cref{rem:g5:human-reproduction:adults}多年前便许诺过本章。现在兑现：
事实、时间表那著名的不公平，
以及机制——它将带来一个比\omterm{def:g8:puberty:puberty}{青春期}本身更大的发现：
人体拥有
第二套通讯系统。

\section{发生了什么变化}

\begin{definition}[青春期]\label{def:g8:puberty:puberty}
\emph{青春期}\index{青春期}是身体从儿童变为成人的阶段——
尤其是成人的含义，
即\cref{ch:g8:sexual-reproduction}中的含义：
生殖器官成熟，开始产生
\emph{\omterm{def:g8:sexual-reproduction:gametes}{配子}}，身体也呈现成人的形态。
它是青少年阶段的生物学核心
（\cref{def:g3:stages-of-life:stages}）。
\end{definition}

\begin{proposition}[可见的变化]\label{prop:g8:puberty:changes}
与生长高峰同时，在数年间逐渐出现：
\begin{itemize}
  \item 女孩通常在约九岁至十三岁之间：
        乳房发育，臀部变宽，出现体毛——
        并且\emph{首次月经期}标志着
        \cref{ch:g8:reproductive-systems}中的周期
        开始启动；
  \item 男孩通常晚一两年：嗓音变低（“变声”），
        肩膀变宽，身体和脸部长出体毛，
        睾丸开始产生
        \omterm{def:g4:animal-reproduction:sexual}{精子}；
  \item 两者皆然：\omterm{def:g1:the-five-senses:sense}{皮肤}变得更油（即长痘的年纪），
        \omterm{rem:g7:removing-wastes:sweat}{汗液}发生变化，情绪随身体重建而起伏——
        \omterm{def:g5:brain-in-command:system}{大脑}也在重新布线，朝成人形态发展。
\end{itemize}
\end{proposition}
\admitted

\begin{remark}[时间表因人而异]\label{rem:g8:puberty:timetable}
\Cref{rem:g1:my-body:different}曾就身高说过这一点；
\omterm{def:g8:puberty:puberty}{青春期}则更加明显地重复这一点：同龄人
在这条路上可能相差数年——早熟和晚熟都正常，
变化的先后次序各不相同，而几乎所有人
最终都会到达同一个地方。上述范围很宽，
其中最宽的部分在于，它们对任何一个人
都说明不了多少。
\end{remark}

\section{机制：第二套信使系统}

\begin{definition}[激素初识]\label{def:g8:puberty:hormones}
\omterm{def:g8:puberty:puberty}{青春期}并非由它所改变的器官决定，它是
\emph{受调控}的——由称为
\emph{激素}\index{激素}的化学信使发出指令：它们释放进入
\emph{\omterm{prop:g4:journey-of-food:absorb}{血液}}，由特化的器官分泌，
被送往全身，由能接收它们的器官执行。
\omterm{def:g8:puberty:puberty}{青春期}时，脑底部的一个控制中心被唤醒，
它的激素指挥生殖器官，而生殖器官
\emph{自身}分泌的激素——由
\cref{prop:g4:heart-and-blood:carries}所说的运输系统送往全身——
统领整个重塑过程：生长、毛发、嗓音和周期。
\end{definition}

\begin{example}[一个信号，许多施工地点]\label{ex:g8:puberty:sites}
\omterm{prop:g4:journey-of-food:absorb}{血液}将\omterm{def:g8:puberty:hormones}{激素}同时送往各处——这正是
\omterm{def:g8:puberty:puberty}{青春期}重塑为何是\emph{协调}进行的：嗓音、肩膀和
毛发都跟随同一信使，按同一时间表
行动，尽管这些部位相隔半个身体。没有
\omterm{def:g5:brain-in-command:system}{神经}线路把喉与胡须相连，而是
由血流的广播完成这一点
（\cref{ch:g8:hormonal-communication}专门讲解这一系统，
为此专设一章）。
\end{example}

\begin{remark}[为何会有情绪的“天气”？]\label{rem:g8:puberty:moods}
同样的信使也到达\omterm{def:g5:brain-in-command:system}{大脑}，而\omterm{def:g5:brain-in-command:system}{大脑}本身
正在改造：情绪起伏、强度变化，\omterm{prop:g1:healthy-habits:needs}{睡眠}
逐渐推迟——这是天气，而非过错。了解这一
机制对双方都有帮助：情绪的风暴是化学与
重建共同作用的结果，它们是正常的，
并且终会过去。
\end{remark}

\section{在重建中生活}

\begin{proposition}[重建期间的照护]\label{prop:g8:puberty:care}
处于重建中的身体有建造者般的需求，这些在本书
前文中都已提到：
\begin{itemize}
  \item 物质与燃料：生长高峰以物质为代价——
        均衡饮食
        （\cref{met:g3:balanced-diet:check}），
        包括构建物质
        （\cref{exo:g7:digestion-and-nutrients:11}）；
  \item \omterm{prop:g1:healthy-habits:needs}{睡眠}：\omterm{def:g5:brain-in-command:system}{大脑}的更新与身体的
        生长都在夜间进行
        （\cref{prop:g5:brain-in-command:protect}）——而
        推迟了的\omterm{def:g6:exploring-our-environment:components}{生物}钟仍需要它应有的\omterm{prop:g1:healthy-habits:needs}{睡眠}时长；
  \item 运动：骨与\omterm{def:g4:heart-and-blood:heart}{心脏}正处于它们
        最后也是最好的建造期
        （\cref{met:g3:bones-and-muscles:care}）；
  \item 并且不得有破坏：
        \cref{prop:g5:healthy-choices:substances}所述的物质
        对仍在建造中的身体和\omterm{def:g5:brain-in-command:system}{大脑}
        伤害最大。
\end{itemize}
\end{proposition}
\admitted

\begin{example}[问题有去处]\label{ex:g8:puberty:questions}
\omterm{def:g8:puberty:puberty}{青春期}会引出许多问题——关于时机、关于身体、
关于感受——这些问题都有合适的去处：学校
护士和医生每天为其解答，且保密，
并不大惊小怪；值得信赖的成年人都曾经历过；
本书各章则便于反复阅读。
不值得信任的是：操场上的传说，以及
网络上最糟糕的角落，它们正是靠
这一年龄段的不确定性牟利。
\end{example}

\section{练习}

\begin{exercise}[$\star$]\label{exo:g8:puberty:1}
给出\omterm{def:g8:puberty:puberty}{青春期}的定义，并指出\omterm{def:g8:puberty:puberty}{青春期}是
\cref{def:g3:stages-of-life:stages}中哪一阶段的核心。
\end{exercise}

\begin{exercise}[$\star$]\label{exo:g8:puberty:2}
列出女孩\omterm{def:g8:puberty:puberty}{青春期}的三种变化、男孩的三种变化，
以及两者共有的两种变化。
\end{exercise}

\begin{exercise}[$\star$]\label{exo:g8:puberty:3}
男女两性中，\omterm{def:g8:sexual-reproduction:gametes}{配子}产生的开始各以什么为标志，
根据\cref{prop:g8:puberty:changes}？
\end{exercise}

\begin{exercise}[$\star$]\label{exo:g8:puberty:4}
什么是\omterm{def:g8:puberty:hormones}{激素}——在何处释放、以何种方式运输、
由谁来响应？
\end{exercise}

\begin{exercise}[$\star$]\label{exo:g8:puberty:5}
\omterm{def:g8:puberty:puberty}{青春期}的指令链从何处开始，
到何处结束？
\end{exercise}

\begin{exercise}[$\star$]\label{exo:g8:puberty:6}
两名十三岁的学生处于明显不同的
阶段。\cref{rem:g8:puberty:timetable}是怎么说的
——请完整说明？
\end{exercise}

\begin{exercise}[$\star\star$]\label{exo:g8:puberty:7}
为何一种信使能够协调相距半个身体的建造部位，
而\omterm{def:g5:brain-in-command:system}{神经}需要分别接线到每一处？请结合
\cref{ex:g8:puberty:sites}作答。
\end{exercise}

\begin{exercise}[$\star\star$]\label{exo:g8:puberty:8}
解释\cref{rem:g8:puberty:moods}中的情绪波动，
——两个原因，一条宽慰。
\end{exercise}

\begin{exercise}[$\star\star$]\label{exo:g8:puberty:9}
把生长突增与健康餐盘联系起来：重建过程
最需要哪几类\omterm{def:g7:digestion-and-nutrients:families}{营养物质}？
其依据来自哪一章的推理？
\end{exercise}

\begin{exercise}[$\star\star$]\label{exo:g8:puberty:10}
为何\cref{prop:g5:healthy-choices:substances}所述的
物质对青少年身体的伤害最大？（“建造”
这一比喻已经给出答案。）
\end{exercise}

\begin{exercise}[$\star\star$]\label{exo:g8:puberty:11}
按可信与否分类，并各给出一条理由：
学校护士；一段病毒式视频中的奇迹说法；本书；
操场上的传说；家庭医生。
\end{exercise}

\begin{exercise}[$\star\star\star$]\label{exo:g8:puberty:12}
“\omterm{def:g8:puberty:puberty}{青春期}是\omterm{def:g8:puberty:hormones}{激素}系统的公开亮相。”
论证这句话：重建过程说明了该系统的
作用范围、协调方式，以及它与
\omterm{def:g5:brain-in-command:system}{神经}信号传递的不同——三条论断，三条
观察结果。
\end{exercise}

\section{问题：生长曲线图讨论}

\begin{problem}[{周末问题——四份（虚构）
生长曲线图，结合全章阅读}]\label{pb:g8:puberty:1}
一位校医审阅四张身高随年龄变化的图表，
每张都标有测量日期：L 的曲线在
十一至十三岁间陡升，现在十五岁时趋于平缓；
M 在十三岁时仍缓慢攀升，突增尚未开始；
N 在十四岁时处于突增中段；O 的陡峭段
位于九至十二岁之间，现已进入平台期。

\medskip\noindent\textbf{第一部分——读懂曲线。}
\begin{enumerate}
  \item 将每张图表对应到
        \cref{rem:g3:stages-of-life:bursts}的示意图上：哪一次
        突增正在显现，每位学生处于其中的
        哪个位置？
  \item M 担心自己“落后”了。请起草医生的
        回答，依据
        \cref{rem:g8:puberty:timetable}，结合
        图表平缓攀升所支持的预测。
  \item O 问平台期是否意味着“十二岁时已
        停止生长”。请纠正这种理解：平台期
        结束的是什么，之后又会继续什么？
  \item 医生为何读取\emph{斜率}，而非身高，
        来判断每位学生所处的阶段？
\end{enumerate}

\medskip\noindent\textbf{第二部分——曲线背后。}
\begin{enumerate}[resume]
  \item 说出引发每次突增的机制，以及
        这些指令所经的路径。
  \item 医生在每次就诊时也会追踪
        \omterm{prop:g1:healthy-habits:needs}{睡眠}与饮食，请将二者分别对应到
        \cref{prop:g8:puberty:care}中的相应条目。
  \item N 反映情绪风暴，以及入睡时间
        逐渐推迟。请依据
        \cref{rem:g8:puberty:moods}给出双重解释，
        并说明这对 N 的上学周\omterm{prop:g1:healthy-habits:needs}{睡眠}有何实际影响。
  \item L 私下向医生询问一项图表无法显示的变化。
        请说明，在不虚构提问内容的情况下，
        为何诊室是
        合适的去处（见\cref{ex:g8:puberty:questions}）。
\end{enumerate}

\medskip\noindent\textbf{第三部分——讨论
纪要。}
\begin{enumerate}[resume]
  \item 为每张图表写一句记录：阶段、
        是否正常，以及需要注意的事项。
  \item 四张图表所示年龄相差数年，但医生
        将四张都归入“正常”。请用一句话
        说明理由。
  \item 一位家长请求“想办法让 M 快起来”。
        请结合时间表的说明与
        相关机制作答：有何需要修复之处？
  \item 以本章的一句话总结收尾，供班级公告栏使用：
        \omterm{def:g8:puberty:puberty}{青春期}是什么，
        由谁来调控，以及它需要什么。
\end{enumerate}
\end{problem}
